\chapter*{แผนบริหารการสอนประจำวิชา}
\addcontentsline{toc}{chapter}{แผนบริหารการสอนประจำวิชา}

\begin{tcolorbox}[enhanced, colback=white, colframe=darkNavy, colbacktitle=darkNavy,
    coltitle=white, fonttitle=\bfseries\normalsize,
    title={ข้อมูลทั่วไปของรายวิชา},
    arc=2mm, boxrule=0.8pt, leftrule=4mm,
    top=8pt, bottom=8pt, left=10pt, right=10pt]
\begin{tabularx}{\textwidth}{l>{\raggedright\arraybackslash}X}
\textbf{รหัสวิชา} & PHY4308 \\
\textbf{ชื่อรายวิชา} & OpenXR และการประมวลผลเชิงพื้นที่สู่นวัตกรรมห้องปฏิบัติการสะเต็มอัจฉริยะ \\
 & (OpenXR and Spatial Computing for Next-Generation STEM Laboratories) \\
\textbf{จำนวนหน่วยกิต} & 3(2-2-5) หน่วยกิต (บรรยาย 2 ชั่วโมง ปฏิบัติการ 2 ชั่วโมง ศึกษาด้วยตนเอง 5 ชั่วโมง) \\
\textbf{หลักสูตร} & วิทยาศาสตรบัณฑิต สาขาวิชาฟิสิกส์ประยุกต์ / นวัตกรรมเทคโนโลยีดิจิทัล \\
\textbf{อาจารย์ผู้สอน} & ผู้ช่วยศาสตราจารย์ ดร.ชีวะ ทัศนา \\
\textbf{ภาคการศึกษา/ชั้นปี} & ภาคการศึกษาที่ 1 ชั้นปีที่ 3 หรือ 4 \\
\end{tabularx}
\end{tcolorbox}

\begin{tcolorbox}[enhanced, colback=openxrSoftBg, colframe=openxrTeal, colbacktitle=openxrTeal,
    coltitle=white, fonttitle=\bfseries\normalsize,
    title={คำอธิบายรายวิชา (Course Description)},
    arc=2mm, boxrule=0.8pt, leftrule=4mm,
    top=8pt, bottom=8pt, left=10pt, right=10pt, before skip=8pt]
ศึกษาหลักการและสถาปัตยกรรมของมาตรฐานเปิด OpenXR ในการพัฒนาสภาพแวดล้อมเสมือนจริงข้ามแพลตฟอร์ม ระบบพิกัดสามมิติและปริภูมิอ้างอิงเชิงพื้นที่ การติดตามตำแหน่งและการหมุนแบบ 6 องศาอิสระ (6DoF) เทคนิคการตรวจจับและติดตามโครงกระดูกมือ 21 จุดร่วม (21-Joint Skeletal Hand Tracking) สำหรับการปฏิสัมพันธ์แบบไร้สัมผัส กลไกการคำนวณของเอนจินฟิสิกส์สำหรับการจำลองปรากฏการณ์ทางสะเต็ม (จลนศาสตร์ พลศาสตร์ การชน และสนามแรง) การเขียนโปรแกรมทางวิทยาศาสตร์ด้วยภาษา Python และ PyOpenXR การบูรณาการโมเดลปัญญาประดิษฐ์และคอมพิวเตอร์วิทัศน์ในการอนุมานท่าทาง การออกแบบและพัฒนาห้องปฏิบัติการวิทยาศาสตร์เสมือนจริง และการประเมินผลการเรียนรู้เชิงรุก
\end{tcolorbox}

\begin{tcolorbox}[enhanced, colback=white, colframe=spatialBlue, colbacktitle=spatialBlue,
    coltitle=white, fonttitle=\bfseries\normalsize,
    title={วัตถุประสงค์การเรียนรู้ตามแนวทาง OBE},
    arc=2mm, boxrule=0.8pt, leftrule=4mm,
    top=8pt, bottom=8pt, left=10pt, right=10pt, before skip=8pt]
เมื่อสิ้นสุดการเรียนการสอนรายวิชานี้ ผู้เรียนสามารถ
\begin{enumerate}
    \item อธิบายสถาปัตยกรรม วงรอบการทำงาน และกลไกของมาตรฐาน OpenXR ได้อย่างถูกต้องตามหลักวิชาการ
    \item คำนวณและประยุกต์ใช้การแปลงพิกัดสามมิติ ควอเทอร์เนียน และเมทริกซ์การแปลงในปริภูมิ 6DoF ได้
    \item พัฒนาระบบปฏิสัมพันธ์เชิงพื้นที่ด้วยการตรวจจับโครงร่างมือ 21 จุดร่วมเพื่อควบคุมอุปกรณ์ทดลองเสมือนได้
    \item สร้างแบบจำลองปรากฏการณ์ทางฟิสิกส์และสะเต็มศึกษาด้วยเอนจินฟิสิกส์เชิงตัวเลขได้อย่างแม่นยำ
    \item บูรณาการไลบรารี Python, PyOpenXR, และโมเดลปัญญาประดิษฐ์ในการประมวลผลข้อมูลการทดลองได้
    \item ออกแบบ สร้าง และประเมินห้องปฏิบัติการเสมือนจริงด้านสะเต็มศึกษาที่ทำงานข้ามแพลตฟอร์มได้
\end{enumerate}
\end{tcolorbox}

\begin{tcolorbox}[enhanced, colback=white, colframe=rbruGold!80!black, colbacktitle=rbruGold!80!black,
    coltitle=white, fonttitle=\bfseries\normalsize,
    title={โครงสร้างแผนการจัดการเรียนรู้รายสัปดาห์ (16 สัปดาห์)},
    arc=2mm, boxrule=0.8pt, leftrule=4mm,
    top=8pt, bottom=8pt, left=10pt, right=10pt, before skip=8pt]
\begin{tabularx}{\textwidth}{ccX}
\toprule
\textbf{สัปดาห์} & \textbf{บทที่} & \textbf{หัวข้อเนื้อหาและกิจกรรมการเรียนรู้} \\
\midrule
1--2 & 1 & สถาปัตยกรรม OpenXR, วงรอบการทำงาน และความสำคัญต่อสะเต็มศึกษา \\
3--4 & 2 & ปริภูมิพิกัด 3 มิติ, การแปลงทางเรขาคณิต และการติดตามการเคลื่อนที่ 6DoF \\
5--6 & 3 & การปฏิสัมพันธ์เชิงพื้นที่และการตรวจจับโครงร่างมือ 21 จุดร่วม (XR\_EXT\_hand\_tracking) \\
7--8 & 4 & เอนจินฟิสิกส์, การจำลองวัตถุแข็งเกร็ง, การชน และระบบแรงโน้มถ่วง \\
9 & - & สอบวัดผลกลางภาค (Midterm Assessment) \\
10--11 & 5 & การเขียนโปรแกรมวิทยาศาสตร์และสะเต็มด้วย Python และ PyOpenXR \\
12--13 & 6 & การบูรณาการปัญญาประดิษฐ์และคอมพิวเตอร์วิทัศน์สำหรับห้องปฏิบัติการเสมือน \\
14--15 & 7 & การพัฒนาและประยุกต์ใช้ห้องปฏิบัติการเสมือนสะเต็มศึกษาขั้นสูง (Optics, Waves, Circuits) \\
16 & - & นำเสนอโครงงานนวัตกรรมห้องทดลองเสมือนจริง และสอบวัดผลปลายภาค \\
\bottomrule
\end{tabularx}
\end{tcolorbox}

\begin{tcolorbox}[enhanced, colback=white, colframe=darkSlate, colbacktitle=darkSlate,
    coltitle=white, fonttitle=\bfseries\normalsize,
    title={การวัดและประเมินผลการเรียนรู้},
    arc=2mm, boxrule=0.8pt, leftrule=4mm,
    top=8pt, bottom=8pt, left=10pt, right=10pt, before skip=8pt]
\begin{itemize}
    \item การมีส่วนร่วมในชั้นเรียนและการเรียนรู้เชิงรุก (Class Participation \& Discussion) \hfill 10\%
    \item รายงานผลการทดลองและการเขียนโปรแกรมปฏิบัติการ (Laboratory \& Coding Tasks) \hfill 30\%
    \item การสอบวัดผลกลางภาค (Midterm Examination) \hfill 20\%
    \item โครงงานพัฒนานวัตกรรมห้องปฏิบัติการเสมือนจริงสะเต็มศึกษา (Term Project) \hfill 20\%
    \item การสอบวัดผลปลายภาค (Final Examination) \hfill 20\%
\end{itemize}
\end{tcolorbox}
\clearpage
